% Purpose: Derive the homogeneous isotropic elastodynamic wave equation and its
% P- and S-wave modes used in the Introduction.
% Status: Draft; derivation and notation require academic review.
% Source-of-truth: doc/UNITS/front/Introduction.tex, Eq.
% (momentum_conservation) and Eq. (hooke_law).

\section*{Derivation of the Isotropic Elastodynamic Wave Equation}
\addcontentsline{toc}{section}{
  Derivation of the Isotropic Elastodynamic Wave Equation
}
\label{sec:supp_derivation_wave_equation}

Let $\Omega \subseteq \mathbb{R}^{3}$ be the spatial domain while time
$t\in\mathbb{R}$, with indices $i,j,k\in\{1,2,3\}$ following the Einstein
summation convention \ie{repeated indices are summed}. The displacement is
$\mathbf{u}(\mathbf{x},t)$; the density $\rho(\mathbf{x}) > 0$, and the Lamé
parameters $\lambda(\mathbf{x}), \mu(\mathbf{x}) > 0$ are spatially constant
material parameters. Thus, the derivation concerns a homogeneous, isotropic,
linear-elastic medium; it follows the standard continuum-elastodynamics
treatment in \cite{AkiRichards1980}.
\begin{equation}
  \nabla \cdot \boldsymbol{\sigma} + \mathbf{f} = \rho
  \frac{\partial^2 \mathbf{u}}{\partial t^2}
  \label{eq:supp_cauchy_momentum}
\end{equation}

\subsection*{Kinematics and constitutive relation}

The infinitesimal strain, its trace, and isotropic Hooke law are
\begin{equation*}
  \begin{aligned}
    \epsilon_{ij}
    & = \frac{1}{2}\left(\frac{\partial u_i}{\partial x_j}
    + \frac{\partial u_j}{\partial x_i}\right) = \epsilon_{ji},
    & \epsilon_{kk} & = \frac{\partial u_k}{\partial x_k}
    = \nabla \cdot \mathbf{u}, \\
    \sigma_{ij} & = \lambda\delta_{ij}\epsilon_{kk} + 2\mu\epsilon_{ij} \\
    & = \lambda\delta_{ij}\frac{\partial u_k}{\partial x_k}
    + \mu\left(\frac{\partial u_i}{\partial x_j}
    + \frac{\partial u_j}{\partial x_i}\right),
  \end{aligned}
\end{equation*}
where $\delta_{ij}$ is the Kronecker delta.

\subsection*{Equation of motion}

Since $\partial_{x_j}\lambda = \partial_{x_j}\mu = 0$, differentiation with
respect to $x_j$ gives
\begin{equation}
  \begin{aligned}
    \frac{\partial\sigma_{ij}}{\partial x_j} & = \frac{\partial}{\partial x_j}
    \left[\lambda\delta_{ij}\frac{\partial u_k}{\partial x_k}
      + \mu\left(\frac{\partial u_i}{\partial x_j}
    + \frac{\partial u_j}{\partial x_i}\right)\right] \\
    & = \lambda\frac{\partial}{\partial x_i}
    \left(\frac{\partial u_k}{\partial x_k}\right)
    + \mu\frac{\partial^2u_i}{\partial x_j\partial x_j}
    + \mu\frac{\partial}{\partial x_i}
    \left(\frac{\partial u_j}{\partial x_j}\right) \\
    & = (\lambda + \mu)\frac{\partial}{\partial x_i}
    \left(\frac{\partial u_k}{\partial x_k}\right) + \mu\nabla^2u_i.
  \end{aligned}
  \label{eq:supp_stress_divergence}
\end{equation}
The interchange
$\partial_{x_j}\partial_{x_i}u_j = \partial_{x_i}\partial_{x_j}u_j$ used in
the second line is valid because the displacement components are assumed to be
twice continuously differentiable, $u_j\in C^2(\Omega \times \mathbb{R})$.
Spatially constant $\lambda$ and $\mu$ separately justify taking those
coefficients outside the spatial derivatives.

Newton's second law applied to a differential volume gives a force balance per
unit volume. Its inertial term is mass density times acceleration,
$\rho\partial_t^2u_i$. The stress-divergence term
$\partial_{x_j}\sigma_{ij}$ is the $i$th component of the net internal force
transmitted across the volume boundary. The remaining term $f_i$ is an
externally prescribed body force per unit volume, for example an equivalent
force representation of a seismic source. Thus, for each $i\in\{1,2,3\}$,
conservation of linear momentum gives
\begin{equation}
  \begin{aligned}
    \rho\frac{\partial^2u_i}{\partial t^2}
    & = \frac{\partial\sigma_{ij}}{\partial x_j} + f_i, \\
    \rho\frac{\partial^2\mathbf{u}}{\partial t^2}
    & = (\lambda + \mu)\nabla(\nabla \cdot \mathbf{u})
    + \mu\nabla^2\mathbf{u} + \mathbf{f}.
  \end{aligned}
  \label{eq:supp_navier_cauchy}
\end{equation}
The first line is the component-wise balance. Substituting
\eqref{eq:supp_stress_divergence} into it and collecting the three
components produces the second line, the vector Navier--Cauchy equation in
\eqref{eq:navier_cauchy} of the main text.

\subsection*{P- and S-wave projections}

For source-free body-wave modes, set $\mathbf{f} = \mathbf{0}$. In an unbounded
homogeneous medium \ie{one with no physical boundaries and with a smooth
displacement field that decays sufficiently rapidly at infinity}, the Helmholtz
decomposition expresses $\mathbf{u}(\mathbf{x},t)$ through a scalar potential
$\phi$ and a vector potential $\boldsymbol{\psi}$ \cite{MorseFeshbach1953}.
The Coulomb gauge $\nabla \cdot \boldsymbol{\psi} = 0$, together with the
stated decay condition, excludes nonzero harmonic residuals:
\begin{equation*}
  \begin{aligned}
    \mathbf{u} = \mathbf{u}_\mathrm{P} + \mathbf{u}_\mathrm{S}
    & = \nabla\phi + \nabla \times \boldsymbol{\psi},
    &\nabla \cdot \boldsymbol{\psi}& = 0, \\
    \nabla \cdot \mathbf{u}& = \nabla^2\phi,
    &\nabla \times \mathbf{u}& = -\nabla^2\boldsymbol{\psi}.
  \end{aligned}
\end{equation*}
Taking the divergence of \eqref{eq:supp_navier_cauchy} gives
\begin{equation*}
  \begin{aligned}
    \rho\frac{\partial^2}{\partial t^2}(\nabla \cdot \mathbf{u})
    & = (\lambda + \mu)\nabla^2(\nabla \cdot \mathbf{u})
    +\mu\nabla^2(\nabla \cdot \mathbf{u}) \\
    & = (\lambda + 2\mu)\nabla^2(\nabla \cdot \mathbf{u}), \\
    \frac{\partial^2}{\partial t^2}(\nabla^2\phi)
    & = \frac{\lambda + 2\mu}{\rho}\nabla^2(\nabla^2\phi).
  \end{aligned}
\end{equation*}
Equivalently,
\begin{equation*}
  \nabla^2\left[
    \frac{\partial^2\phi}{\partial t^2} -
    \frac{\lambda + 2\mu}{\rho}\nabla^2\phi
  \right] = 0.
\end{equation*}
The bracketed residual is harmonic. The gauge and decay conditions exclude a
nonzero harmonic residual, so the scalar potential satisfies
\begin{equation}
  \begin{aligned}
    \frac{\partial^2\phi}{\partial t^2} & = \alpha^2\nabla^2\phi,
    &\alpha^2& = \frac{\lambda + 2\mu}{\rho}.
  \end{aligned}
  \label{eq:supp_p_wave_equation}
\end{equation}

Taking the curl of \eqref{eq:supp_navier_cauchy} gives
\begin{equation*}
  \begin{aligned}
    \rho\nabla \times \frac{\partial^2\mathbf{u}}{\partial t^2}
    & = (\lambda + \mu)\nabla \times \nabla(\nabla \cdot \mathbf{u})
    +\mu\nabla \times \nabla^2\mathbf{u} \\
    & = \mathbf{0} + \mu\nabla^2(\nabla \times \mathbf{u}) \\
    & = \mu\nabla^2(\nabla \times \mathbf{u}).
  \end{aligned}
\end{equation*}
The first term vanishes because the curl of a gradient is zero \ie{
$\nabla \times \nabla f = 0$ for any smooth scalar field $f$}. The temporal
derivatives commute with curl, and curl commutes with the Laplacian because the
field is sufficiently smooth. Therefore,
\begin{equation*}
  \begin{aligned}
    \rho\frac{\partial^2}{\partial t^2}(\nabla \times \mathbf{u})
    & = \mu\nabla^2(\nabla \times \mathbf{u}), \\
    \rho\frac{\partial^2}{\partial t^2}(-\nabla^2\boldsymbol{\psi})
    & = \mu\nabla^2(-\nabla^2\boldsymbol{\psi}), \\
    - \rho\nabla^2\frac{\partial^2\boldsymbol{\psi}}{\partial t^2}
    & = -\mu\nabla^4\boldsymbol{\psi}.
  \end{aligned}
\end{equation*}
Equivalently,
\begin{equation*}
  \nabla^2\left[\frac{\partial^2\boldsymbol{\psi}}{\partial t^2}
    - \frac{\mu}{\rho}\nabla^2\boldsymbol{\psi}
  \right] = \mathbf{0}.
\end{equation*}
The bracketed vector residual is harmonic. Under the same gauge and decay
conditions, it vanishes, so the shear potential satisfies
\begin{equation}
  \begin{aligned}
    \frac{\partial^2\boldsymbol{\psi}}{\partial t^2}
    & = \beta^2\nabla^2\boldsymbol{\psi},
    &\beta^2& = \frac{\mu}{\rho}.
  \end{aligned}
  \label{eq:supp_s_wave_equation}
\end{equation}
Thus the P mode depends on the bulk-plus-shear stiffness $\lambda + 2\mu$
whereas the S mode only depends on the shear modulus $\mu$. Under the
main text's conventional conditions $\rho > 0$, $\lambda > 0$, and $\mu > 0$,
\begin{equation*}
  \alpha^2 - \beta^2 = \frac{\lambda + \mu}{\rho} > 0
  \quad\Longrightarrow\quad \alpha > \beta.
\end{equation*}

\subsection*{Alternative curl--curl formulation}

The preceding divergence and curl projections establish the scalar and vector
potential equations directly. The equivalent curl--curl formulation below
reconstructs the same separation at the displacement level, making explicit how
volumetric and rotational motions enter the source-free Navier--Cauchy equation
under the same gauge and decay conditions.

Using the vector identity for the Laplacian of a vector field,
\begin{equation*}
  \nabla^2\mathbf{u}
  = \nabla(\nabla \cdot \mathbf{u}) - \nabla \times (\nabla \times \mathbf{u})
\end{equation*}
the source-free Navier--Cauchy equation \eqref{eq:supp_navier_cauchy} becomes
\begin{equation}
  \begin{aligned}
    \rho\frac{\partial^2\mathbf{u}}{\partial t^2}
    & = (\lambda + \mu)\nabla(\nabla \cdot \mathbf{u})
    + \mu\nabla^2\mathbf{u} \\
    & = (\lambda + \mu)\nabla(\nabla \cdot \mathbf{u})
    + \mu\left[\nabla(\nabla \cdot \mathbf{u})
    - \nabla \times (\nabla \times \mathbf{u})\right] \\
    & = (\lambda + 2\mu)\nabla(\nabla \cdot \mathbf{u})
    - \mu\nabla \times (\nabla \times \mathbf{u}).
  \end{aligned}
  \label{eq:supp_curl_curl_navier_cauchy}
\end{equation}
The first term is the irrotational force associated with volumetric dilation or
compression. The second is the rotational force associated with equivoluminal
shear deformation.

Substitution into \eqref{eq:supp_curl_curl_navier_cauchy} gives
\begin{equation*}
  \begin{aligned}
    \rho\frac{\partial^2}{\partial t^2}
    \left(\nabla\phi + \nabla \times \boldsymbol{\psi}\right)
    & = (\lambda + 2\mu)\nabla(\nabla^2\phi) - \mu\nabla \times \nabla \times
    \left(\nabla \times \boldsymbol{\psi}\right), \\
    \nabla \times \nabla \times \boldsymbol{\psi}
    & = \nabla(\nabla \cdot \boldsymbol{\psi}) - \nabla^2\boldsymbol{\psi}
    = -\nabla^2\boldsymbol{\psi}, \\
    \nabla \times \nabla \times \left(\nabla \times \boldsymbol{\psi}\right)
    & = -\nabla^2\left(\nabla \times \boldsymbol{\psi}\right).
  \end{aligned}
\end{equation*}
Consequently, the irrotational and solenoidal contributions satisfy
\begin{equation}
  \nabla\left[\rho\frac{\partial^2\phi}{\partial t^2}
  - (\lambda + 2\mu)\nabla^2\phi\right]
  + \nabla \times \left[\rho\frac{\partial^2\boldsymbol{\psi}}{\partial t^2}
  - \mu\nabla^2\boldsymbol{\psi}\right] = \mathbf{0}.
  \label{eq:supp_helmholtz_decoupling}
\end{equation}
Taking the divergence of \eqref{eq:supp_helmholtz_decoupling} gives the
harmonic scalar residual derived above. Taking its curl gives a vanishing
double curl of the vector residual; the Coulomb gauge makes that residual
divergence-free, so the vector identity implies that it is harmonic as well.
Under the stated decay condition, both residuals vanish. Therefore,
\eqref{eq:supp_helmholtz_decoupling} recovers the previously obtained
independent potential equations:
\begin{equation}
  \rho\frac{\partial^2\phi}{\partial t^2} = (\lambda + 2\mu)\nabla^2\phi,\qquad
  \rho\frac{\partial^2\boldsymbol{\psi}}{\partial t^2} =
  \mu\nabla^2\boldsymbol{\psi}.
\end{equation}

\subsection*{Plane-wave check}

Let $\mathbf{k}_\phi, \mathbf{k}_\psi \in \mathbb{R}^{3}$ be the wave vectors,
the angular frequencies be $\omega_\phi > 0$ and $\omega_\psi > 0$, the scalar
amplitude be $A\in\mathbb{C}$, and $\mathbf{B}\in\mathbb{C}^{3}$ the vector
amplitude with $\mathbf{k}_\psi \cdot \mathbf{B} = 0$. For
\begin{equation*}
  \begin{aligned}
    \phi & = A\exp\!\left(\mathrm{i}(\mathbf{k}_\phi \cdot \mathbf{x} -
    \omega_\phi t)\right),
    &\boldsymbol{\psi} & = \mathbf{B}\exp\!\left(
    \mathrm{i}(\mathbf{k}_\psi \cdot \mathbf{x} - \omega_\psi t)\right), \\
    \frac{\partial^2\phi}{\partial t^2} & = -\omega_\phi^2\phi,
    &\nabla^2\phi & = -\lVert\mathbf{k}_\phi\rVert^2\phi, \\
    \frac{\partial^2\boldsymbol{\psi}}{\partial t^2}
    & = -\omega_\psi^2\boldsymbol{\psi},
    &\nabla^2\boldsymbol{\psi}
    & = -\lVert\mathbf{k}_\psi\rVert^2\boldsymbol{\psi}.
  \end{aligned}
\end{equation*}
Substitution into the two potential equations (\eqref{eq:supp_p_wave_equation},
\eqref{eq:supp_s_wave_equation}) gives
\begin{equation}
  \omega_\phi^2 = \alpha^2\lVert\mathbf{k}_\phi\rVert^2,
  \qquad v_{\mathrm{P}} =
  \frac{\omega_\phi}{\lVert\mathbf{k}_\phi\rVert} = \alpha.
  \label{eq:supp_p_dispersion}
\end{equation}
The same substitution for $\boldsymbol{\psi}$ gives
\begin{equation}
  \omega_\psi^2 = \beta^2\lVert\mathbf{k}_\psi\rVert^2,
  \qquad v_{\mathrm{S}} =
  \frac{\omega_\psi}{\lVert\mathbf{k}_\psi\rVert} = \beta.
  \label{eq:supp_s_dispersion}
\end{equation}
These relations reproduce the P- and S-wave velocities stated in
\eqref{eq:wave_velocities}. The result is limited to the homogeneous,
isotropic, linear-elastic, source-free body-wave derivation; interfaces,
attenuation, anisotropy, and free-surface effects require additional boundary
conditions and constitutive terms.
